% Page 060 — page imprimée 57
% Le Pasteur Etienne Albert Giran

\begin{center}
{\Large\bfseries Le Pasteur Etienne Albert Giran}
\end{center}

\begin{itemize}
\item Entre les deux guerres, rarissimes étaient les familles qui venaient passer des vacances
à Meyrueis, dans une maison familiale, sans en être originaires. Il y avait bien
quelques rares locations, mais les touristes venant « faire les Gorges du Tarn », visiter
Dargilan ou s’aventurer à Bramabiau prenaient pension dans les hôtels ou, pour les
plus fortunés au Château d’Ayres.

\item Or le Pasteur Giran possédait une belle maison, en plein centre ville en face de
l’Eglise et il n’était pas, tout comme sa femme originaire de Meyrueis. Comment se
retrouvaient-ils propriétaires, rue du Champ de Mars, à Meyrueis ?
\end{itemize}

Les familles Giran sont originaires de Nîmes et des environs. Etienne Giran avait fait ses
études de théologie à la Faculté de théologie de Montauban. De tendance libérale, on peut
imaginer qu’il avait été appelé comme pasteur, par les Eglises Wallonnes de Hollande, peut-
être même par son collègue le pasteur Henri Valès originaire de Meyrueis qui avait épousé
une Marie Suzanne Martin elle aussi de Meyrueis. Il avait ses entrées à la cour de la reine de
Hollande Wilhelmine.

Henri Valès devait apprécier son jeune collègue Giran puisque après sa mort, sa veuve, Marie
Suzanne Martin lui légua, en 1903, leur maison de Meyrueis.

C’est ainsi que les Giran passèrent une partie de leurs congés à Meyrueis et qu’ils firent
profiter de cette maison nombre de leurs collègues ou amis, jusqu’à la famille Roth qui fut la
dernière à l’occuper avant sa vente après la guerre 39-45.

\vspace{1em}

« Albert Giran : Pasteur, philosophe et écrivain », écrit Florian Hollard dans un livre consacré
à son père « Michel Hollard, le Français qui a sauvé Londres ».

Théologien libéral il exerça son ministère, avant la guerre 39-45, dans les Eglises Wallonnes
et eut deux enfants Olivier et Claire.

Il était aussi journaliste à l’AFP et il prit sa retraite à Sèvres au début de la guerre. Il inspirait
un grand respect pour avoir eu le courage de blâmer ouvertement la fameuse poignée de main
échangée entre Pétain et Hitler à Montoire. Il était aussi l’auteur d’un pamphlet anti-allemand :
« Devant l’Histoire » paru dans le journal « Réforme » du 14 Juillet 1945.

Son fils \textbf{Olivier}, inspiré par l’exemple de son père voulut très jeune s’engager dans un réseau
de résistance. Mis en relation avec le réseau Agir, il offrit ses services à Michel Hollard qui
devant son enthousiasme, son intelligence, et aussi sa connaissance de certains tronçons de la
frontière franco-suisse lui confia la responsabilité d’acheminement de renseignements aux
services anglais basés en Suisse.

\begin{quote}\itshape
« Il établit et assure, de 1941 à 1942, la liaison avec le Commandement par delà une frontière militairement
gardée par l’ennemi... »
\end{quote}

\noindent\textit{extrait d’une citation qui accompagne sa nomination dans l’Ordre de la Légion d’Honneur.}

Il passe en Suisse tous les quinze jours et y accède par des chemins qu’il connaît bien, car
depuis son enfance il a fait des séjours nombreux à la Côte aux Fées, village proche de la
frontière dans une pension de famille, la Crête, dont il parle plusieurs fois dans ses lettres.

En dehors des missions dont il a pris la charge et qui constituent pour lui le service commandé,
il conduit en Suisse, à deux reprises des Hollandais qui ont de sérieuses raisons d’échapper
aux poursuites de la police allemande.
